\documentclass[compress,handout]{beamer}
\mode<presentation>
{
 \usetheme{Vilanova}
}

\usepackage[french]{babel}

\usepackage[utf8]{inputenc}

\usepackage{times}
\usepackage[T1]{fontenc}

\usepackage{amsfonts}
\usepackage{amsmath}
\usepackage{amssymb}
\usepackage{tikz}
\usepackage{eurosym}
%\usepackage{url}
\usepackage[normal]{subfigure}
\newcommand{\goodgap}{%
	\hspace{\subfigtopskip}%
	\hspace{\subfigbottomskip}}



%\newtheorem{definition}{Definition}

\title{Traitement d'images de télédétection}
\subtitle{Détection de changements}




\author
{jordi.inglada@cesbio.cnes.fr}
\normalsize

\institute[Cesbio] % (optional, but mostly needed)
{\textsc{Centre d'Études Spatiales de la Biosphère, Toulouse, France}}

\date{}

\pgfdeclareimage[height=96mm,width=128mm]{background}{fondsClairSansLogo}
\setbeamertemplate{background}{\pgfuseimage{background}}
\pgfdeclareimage[height=0.6cm]{logoIncrust}{logoIncrust}
\pgfdeclareimage[height=0.5cm]{logo_cesbio}{logo_cesbio}
\logo{
\begin{tabular}{lp{0.25\textwidth}lp{0.25\textwidth}r}
\href{http://www.cesbio.ups-tlse.fr/}{\pgfuseimage{logo_cesbio}}
&&\footnotesize{AUF - Marrakech 2011}&&
\href{http://www.orfeo-toolbox.org}{\pgfuseimage{logoIncrust}}\\
\end{tabular}
}


\subject{Détection de changements}




% Delete this, if you do not want the table of contents to pop up at
% the beginning of each subsection:
\AtBeginSubsection[]
{
  \begin{frame}<beamer>
    \frametitle{Outline}
    \tableofcontents[currentsection,currentsubsection]
  \end{frame}
}




% If you wish to uncover everything in a step-wise fashion, uncomment
% the following command: 

%\beamerdefaultoverlayspecification{<+->}
\begin{document}

\begin{frame}
  \titlepage
  \begin{center}
{\tiny Ce contenu est dérivé de la formation \href{http://www.orfeo-toolbox.org/packages/PragmaticRemoteSensing-handout.pdf}{``Pragmatic Remote
  Sensing''} dispensée par J. Inglada et E. Christophe en juillet 2010
  dans le cadre du colloque IGARSS. Il est mis à disposition selon les termes de la licence :\\
Creative Commons Paternité – Partage à l’Identique 3.0 non transcrit.} \href{http://creativecommons.org/licenses/by-sa/3.0/}{\includegraphics[width=0.05\textwidth]{/home/inglada/Dev/GH/IGARSS2010/Tutorial/Slides/Ressources/CC-licence.png}}    
  \end{center}
\end{frame}


\section[Stratégies]{Stratégies classiques pour la détections de changements}

\begin{frame}
  \frametitle{Approches possibles}
\begin{itemize}

\item{Stratégie $1$ : détecteurs simples}

  Production d'une image de vraisemblance de changement (différences,
  ratios, etc.) et seuillage pour produire une carte binaire.

\item{Stratégie $2$ : Comparaison post-classification}

Génération de 2 cartes d'occupation des sols (une pour chaque date) et
comparaison des classes.

\item{Stratégie $3$ : Classification conjointe}

  Génération de la carte de changements directement à partir de la
  classification conjointe des 2 images.


\end{itemize}
\end{frame}


\section[Détecteurs]{Détecteurs disponibles dans l'OTB}
\begin{frame}
  \frametitle{Détecteurs disponibles}
  \small
  \begin{itemize}
    \item Différence pixel à pixel des valeurs moyennes dans les
      voisinages :
      \begin{equation}
	I_{D}(i,j)=I_{2}(i,j)-I_{1}(i,j).
      \end{equation}
    \item Ratio de moyennes locales :
      \begin{equation}
\displaystyle I_{R}(i,j) = 1 - min \left(\frac{\displaystyle I_{2}(i,j)}{\displaystyle I_{1}(i,j)},\frac{\displaystyle I_{1}(i,j)}{\displaystyle I_{2}(i,j)}\right).
\end{equation}
      \item Corrélation locale :
\begin{equation}
  I_\rho(i,j) = \frac{1}{N}\frac{\sum_{i,j}(I_1(i,j)-m_{I_1})(I_2(i,j)-m_{I_2})}{\sigma_{I_1}
\sigma_{I_2}}
\end{equation}
\item Distance de Kullback-Leibler entre les distributions locales
  \item Plusieurs versions de l'information mutuelle
  \end{itemize}
  \normalsize
\end{frame}

\begin{frame}
  \frametitle{La main à la pâte}
  \framesubtitle{Affichage des différences}
  \begin{enumerate}
  \item Images : SpotAfter et SpotBefore dans le dossier /home/auf/OTB/data/Examples
  \item Monteverdi: File $\rightarrow$ Concatenate Images
  \item Choisir les amplitudes des 2 images et construire 1 image à 2 bandes
  \item Monteverdi: Visualization $\rightarrow$ Viewer
  \item Choisir l'image à 2 bandes
  \item Dans l'onglet {\em Setup}, choisir {\em RGB composition mode}
    et prendre 1,2,2.
  \item Interpréter les couleurs observées
  \item La même chose pourrait être faite en utilisant des images de primitives
  \end{enumerate}
\end{frame}

\begin{frame}
  \frametitle{La main à la pâte}
  \framesubtitle{Seuillage des différences}
  \begin{enumerate}
  \item Monteverdi: Filtering $\rightarrow$ Band Math
  \item Choisir les amplitudes des 2 images et calculer une différence
  \item Monteverdi: Filtering $\rightarrow$ Threshold
  \item Choisir l'image de différences et appliquer différents seuils
  \item La même chose pourrait être faite avec des ratios
  \item La même chose pourrait être faite en utilisant des images de primitives
  \end{enumerate}
\end{frame}


\section[Application]{Détection de changements interactive}

\begin{frame}
  \frametitle{Détection de changements interactive}
  \begin{itemize}
    \item Génération de cartes de changements binaires en utilisant une
      IHM
    \item Utilise des détecteurs simples en entrée
    \item L'opérateur donne des exemples de {\em changement} et {\em non-changement}
    \item Classification supervisée par SVM
  \end{itemize}
\end{frame}


\begin{frame}
  \frametitle{La main à la pâte}
  \framesubtitle{Classification conjointe}
  \begin{enumerate}
  \item Monteverdi: Filtering $\rightarrow$ Change Detection
  \item Choisir les 2 images à traiter
    \begin{itemize}
    \item Images à 2 dates différentes
    \end{itemize}
  \item Décocher {\em Use Change Detectors}
  \item Utiliser les boutons {\em Changed/Unchanged Class} pour
    choisir chacune des classes
  \item Tracer des polygones sur les images pour construire les
    échantillons d'apprentissage
  \item Le bouton {\em End Polygon} sert à fermer les polygones
  \item Après avoir sélectionné plusieurs polygones par classe cliquer
    sur {\em Learn}
  \item Le bouton {\em Display Results} permet d'afficher les résultats
  \end{enumerate}
\end{frame}

\begin{frame}
  \frametitle{La main à la pâte}
  \framesubtitle{Classification conjointe avec détecteurs de changement}
  \begin{enumerate}
  \item Monteverdi: Filtering $\rightarrow$ Change Detection
      \item Choisir les 2 images à traiter
    \begin{itemize}
    \item Images à 2 dates différentes
    \end{itemize}
  \item S'assurer que la case {\em Use Change Detectors} est cochée
  \item Procéder comme dans l'exercice précédent
  \end{enumerate}
\end{frame}

\begin{frame}
  \frametitle{La main à la pâte}
  \framesubtitle{Classification conjointe avec primitives}
  \begin{enumerate}
  \item Monteverdi: Filtering $\rightarrow$ Change Detection
      \item Choisir les 2 images à traiter
    \begin{itemize}
    \item Images de primitives à 2 dates différentes
    \end{itemize}
  \item Procéder comme dans l'exercice précédent
  \end{enumerate}
\end{frame}



\end{document}
